% GENERATED FILE. Edit canonical lessons, then run npm run generate:books.
% canonical-source-hash: fnv1a64:31ef0f3af4025434
% canonical-lessons: KA-C411-sarattu, KA-C411-oppanda, KA-C411-odambadike, KA-C411-iderisu, KA-C411-badhyate, KA-R411-review-405-411

\chapter{Terms and Conditions}
\label{ch:ka-terms-and-conditions}
\hlchaptermodality{411}

\begin{chapteropening}
\textbf{By the end of this chapter:} \emph{I can say a condition, an agreement, a pact, to fulfil and an obligation.}

Complete the last lesson of chapter 411: I can say a condition, an agreement, a pact, to fulfil and an obligation.
\end{chapteropening}

\section[\texorpdfstring{\kn{ಷರತ್ತು} (ṣarattu)}{ಷರತ್ತು (ṣarattu)}]{\kn{ಷರತ್ತು} (ṣarattu) — a condition; a term}

\label{lesson:KA-C411-sarattu}



\begin{quote}
Before the new one: say the Kannada for above all, then the Kannada for not only that.
\end{quote}

\subsection*{You'll want to know: \kn{ಷರತ್ತು}}
\textbf{\kn{ಷರತ್ತು}} — \emph{ṣarattu} — \textquotedblleft{}a condition; a term\textquotedblright{}.

\begin{quote}
\kn{ಒಂದು} \kn{ಷರತ್ತಿನ} \kn{ಮೇಲೆ} \kn{ನಾನು} \kn{ಒಪ್ಪುತ್ತೇನೆ}.
\end{quote}

\begin{quote}
\emph{ondu ṣarattina mēle nānu opputtēne.}
\end{quote}

\begin{quote}
I agree on one condition.
\end{quote}

\kn{ಷರತ್ತಿನ} \kn{ಮೇಲೆ}, on condition. \kn{ಷರತ್ತುಗಳು}, terms.

\begin{cousinweb}
\textbf{\kn{ಷರತ್ತು}} (\emph{ṣarattu}) is a loan from Arabic \emph{sharṭ}, \textquotedblleft{}a condition\textquotedblright{}, which reached Kannada through Hindustani.
\end{cousinweb}

\begin{grammarlens}[title={What you've built}]
One more word for saying what happens if.
\end{grammarlens}

\subsection*{Your turn}
\begin{itemize}
\raggedright
  \item \emph{Say it:} \emph{ṣarattu}
  \item \emph{Say it:} \emph{ṣarattu}, once more
  \item \emph{Say it:} \emph{ondu ṣarattina mēle nānu opputtēne.} — I agree on one condition.
  \item \emph{Recall:} say the Kannada for above all, then the Kannada for not only that, then say \emph{ṣarattu} again
\end{itemize}

\begin{tcolorbox}[breakable,colback=teal!4,colframe=teal!35!black,title={Before you move on}]
What does \kn{ಷರತ್ತು} mean? (\textquotedblleft{}A condition; a term\textquotedblright{}.) Say it once more.
\end{tcolorbox}
\section[\texorpdfstring{\kn{ಒಪ್ಪಂದ} (oppanda)}{ಒಪ್ಪಂದ (oppanda)}]{\kn{ಒಪ್ಪಂದ} (oppanda) — an agreement; a deal}

\label{lesson:KA-C411-oppanda}



\begin{quote}
Before the new one: say the Kannada for not only that, then the Kannada for a condition.
\end{quote}

\subsection*{You'll want to know: \kn{ಒಪ್ಪಂದ}}
\textbf{\kn{ಒಪ್ಪಂದ}} — \emph{oppanda} — \textquotedblleft{}an agreement; a deal\textquotedblright{}.

\begin{quote}
\kn{ಎರಡು} \kn{ಕಂಪನಿಗಳು} \kn{ಒಂದು} \kn{ಒಪ್ಪಂದ} \kn{ಮಾಡಿಕೊಂಡವು}.
\end{quote}

\begin{quote}
\emph{eraḍu kampanigaḷu ondu oppanda māḍikoṇḍavu.}
\end{quote}

\begin{quote}
The two companies made an agreement.
\end{quote}

\kn{ಒಪ್ಪಂದ} \kn{ಮಾಡಿಕೊ}, to make a deal.

\begin{cousinweb}
\textbf{\kn{ಒಪ್ಪಂದ}} (\emph{oppanda}) is native Dravidian: DEDR 924 lists Kannada \emph{oppanda}, \textquotedblleft{}an agreement\textquotedblright{}, beside \emph{oppu}, \textquotedblleft{}to agree\textquotedblright{}.
\end{cousinweb}

\begin{grammarlens}[title={What you've built}]
One more word for saying what happens if.
\end{grammarlens}

\subsection*{Your turn}
\begin{itemize}
\raggedright
  \item \emph{Say it:} \emph{oppanda}
  \item \emph{Say it:} \emph{oppanda}, once more
  \item \emph{Say it:} \emph{eraḍu kampanigaḷu ondu oppanda māḍikoṇḍavu.} — The two companies made an agreement.
  \item \emph{Recall:} say the Kannada for not only that, then the Kannada for a condition, then say \emph{oppanda} again
\end{itemize}

\begin{tcolorbox}[breakable,colback=teal!4,colframe=teal!35!black,title={Before you move on}]
What does \kn{ಒಪ್ಪಂದ} mean? (\textquotedblleft{}An agreement; a deal\textquotedblright{}.) Say it once more.
\end{tcolorbox}
\section[\texorpdfstring{\kn{ಒಡಂಬಡಿಕೆ} (oḍambaḍike)}{ಒಡಂಬಡಿಕೆ (oḍambaḍike)}]{\kn{ಒಡಂಬಡಿಕೆ} (oḍambaḍike) — a pact; a treaty}

\label{lesson:KA-C411-odambadike}



\begin{quote}
Before the new one: say the Kannada for a condition, then the Kannada for an agreement.
\end{quote}

\subsection*{You'll want to know: \kn{ಒಡಂಬಡಿಕೆ}}
\textbf{\kn{ಒಡಂಬಡಿಕೆ}} — \emph{oḍambaḍike} — \textquotedblleft{}a pact; a treaty\textquotedblright{}.

\begin{quote}
\kn{ಎರಡು} \kn{ದೇಶಗಳ} \kn{ನಡುವೆ} \kn{ಹೊಸ} \kn{ಒಡಂಬಡಿಕೆ} \kn{ಆಯಿತು}.
\end{quote}

\begin{quote}
\emph{eraḍu dēśagaḷa naḍuve hosa oḍambaḍike āyitu.}
\end{quote}

\begin{quote}
A new treaty was made between the two countries.
\end{quote}

More formal than \kn{ಒಪ್ಪಂದ}.

\begin{cousinweb}
\textbf{\kn{ಒಡಂಬಡಿಕೆ}} (\emph{oḍambaḍike}) is native Dravidian: DEDR 945 lists Kannada \emph{oḍam-baḍu}, \textquotedblleft{}to agree\textquotedblright{}, and \emph{oḍa-baḍike}, \textquotedblleft{}an agreement, a treaty\textquotedblright{}.
\end{cousinweb}

\begin{grammarlens}[title={What you've built}]
One more word for saying what happens if.
\end{grammarlens}

\subsection*{Your turn}
\begin{itemize}
\raggedright
  \item \emph{Say it:} \emph{oḍambaḍike}
  \item \emph{Say it:} \emph{oḍambaḍike}, once more
  \item \emph{Say it:} \emph{eraḍu dēśagaḷa naḍuve hosa oḍambaḍike āyitu.} — A new treaty was made between the two countries.
  \item \emph{Recall:} say the Kannada for a condition, then the Kannada for an agreement, then say \emph{oḍambaḍike} again
\end{itemize}

\begin{tcolorbox}[breakable,colback=teal!4,colframe=teal!35!black,title={Before you move on}]
What does \kn{ಒಡಂಬಡಿಕೆ} mean? (\textquotedblleft{}A pact; a treaty\textquotedblright{}.) Say it once more.
\end{tcolorbox}
\section[\texorpdfstring{\kn{ಈಡೇರಿಸು} (īḍērisu)}{ಈಡೇರಿಸು (īḍērisu)}]{\kn{ಈಡೇರಿಸು} (īḍērisu) — to fulfil}

\label{lesson:KA-C411-iderisu}



\begin{quote}
Before the new one: say the Kannada for an agreement, then the Kannada for a pact.
\end{quote}

\subsection*{You'll want to know: \kn{ಈಡೇರಿಸು}}
\textbf{\kn{ಈಡೇರಿಸು}} — \emph{īḍērisu} — \textquotedblleft{}to fulfil\textquotedblright{}.

\begin{quote}
\kn{ಅಪ್ಪ} \kn{ನನ್ನ} \kn{ಆಸೆಯನ್ನು} \kn{ಈಡೇರಿಸಿದರು}.
\end{quote}

\begin{quote}
\emph{appa nanna āseyannu īḍērisidaru.}
\end{quote}

\begin{quote}
Father fulfilled my wish.
\end{quote}

Past (I …): \kn{ಈಡೇರಿಸಿದೆ}, \emph{īḍēriside}.

\kn{ಈಡೇರು}, to come true: \kn{ಆಸೆ} \kn{ಈಡೇರಿತು}.

\begin{cousinweb}
\textbf{\kn{ಈಡೇರಿಸು}} (\emph{īḍērisu}) is built as the causative of \emph{īḍēru}, \textquotedblleft{}to be fulfilled\textquotedblright{}; its origin is not settled.
\end{cousinweb}

\begin{grammarlens}[title={What you've built}]
One more word for saying what happens if.
\end{grammarlens}

\subsection*{Your turn}
\begin{itemize}
\raggedright
  \item \emph{Say it:} \emph{īḍērisu}
  \item \emph{Say it:} \emph{īḍērisu}, once more
  \item \emph{Say it:} \emph{appa nanna āseyannu īḍērisidaru.} — Father fulfilled my wish.
  \item \emph{Recall:} say the Kannada for an agreement, then the Kannada for a pact, then say \emph{īḍērisu} again
\end{itemize}

\begin{tcolorbox}[breakable,colback=teal!4,colframe=teal!35!black,title={Before you move on}]
What does \kn{ಈಡೇರಿಸು} mean? (\textquotedblleft{}To fulfil\textquotedblright{}.) Say it once more.
\end{tcolorbox}
\section[\texorpdfstring{\kn{ಬಾಧ್ಯತೆ} (bādhyate)}{ಬಾಧ್ಯತೆ (bādhyate)}]{\kn{ಬಾಧ್ಯತೆ} (bādhyate) — an obligation}

\label{lesson:KA-C411-badhyate}



\begin{quote}
Before the new one: say the Kannada for a pact, then the Kannada for to fulfil.
\end{quote}

\subsection*{You'll want to know: \kn{ಬಾಧ್ಯತೆ}}
\textbf{\kn{ಬಾಧ್ಯತೆ}} — \emph{bādhyate} — \textquotedblleft{}an obligation\textquotedblright{}.

\begin{quote}
\kn{ಸಾಲ} \kn{ತೀರಿಸುವುದು} \kn{ನನ್ನ} \kn{ಬಾಧ್ಯತೆ}.
\end{quote}

\begin{quote}
\emph{sāla tīrisuvudu nanna bādhyate.}
\end{quote}

\begin{quote}
Paying off the loan is my obligation.
\end{quote}

\kn{ಬಾಧ್ಯತೆ} \kn{ಇರು}, to be bound to do.

\begin{cousinweb}
\textbf{\kn{ಬಾಧ್ಯತೆ}} (\emph{bādhyate}) is built on Sanskrit \emph{bādhya}, \textquotedblleft{}to be pressed hard\textquotedblright{} (Monier-Williams), with -\emph{tā}, \textquotedblleft{}-ness\textquotedblright{}: what presses on one.
\end{cousinweb}

\begin{grammarlens}[title={What you've built}]
One more word for saying what happens if.
\end{grammarlens}

\subsection*{Your turn}
\begin{itemize}
\raggedright
  \item \emph{Say it:} \emph{bādhyate}
  \item \emph{Say it:} \emph{bādhyate}, once more
  \item \emph{Say it:} \emph{sāla tīrisuvudu nanna bādhyate.} — Paying off the loan is my obligation.
  \item \emph{Recall:} say the Kannada for a pact, then the Kannada for to fulfil, then say \emph{bādhyate} again
\end{itemize}

\begin{tcolorbox}[breakable,colback=teal!4,colframe=teal!35!black,title={Before you move on}]
What does \kn{ಬಾಧ್ಯತೆ} mean? (\textquotedblleft{}An obligation\textquotedblright{}.) Say it once more.
\end{tcolorbox}
\section[(dialogue)]{Looking back: from Settling Disputes to Terms and Conditions}

\label{lesson:KA-R411-review-405-411}



\begin{quote}
Say the Kannada for an alliance and a negotiation.
\end{quote}

\subsection*{Your turn}
Say these aloud:

\begin{itemize}
\raggedright
  \item an alliance, a negotiation, mediation, a compromise, to incite
  \item a university, a degree, a lecture, a professor, to instruct
  \item research, an experiment, a theory, to analyse, a survey
  \item to focus, a study, optional, a syllabus, revision
  \item an essay, to summarise, a treatise, to present, a note
  \item thus, however, on account of, above all, not only that
  \item a condition, an agreement, a pact, to fulfil, an obligation
\end{itemize}

\begin{tcolorbox}[breakable,colback=teal!4,colframe=teal!35!black,title={Before you move on}]
An alliance? (\textbf{\kn{ಮೈತ್ರಿ}}, \emph{maitri}.) A negotiation? (\textbf{\kn{ಸಂಧಾನ}}, \emph{sandhāna}.) Mediation? (\textbf{\kn{ಮಧ್ಯಸ್ಥಿಕೆ}}, \emph{madhyasthike}.) A compromise? (\textbf{\kn{ಹೊಂದಾಣಿಕೆ}}, \emph{hondāṇike}.) To incite? (\textbf{\kn{ಪ್ರಚೋದಿಸು}}, \emph{pracōdisu}.) A university? (\textbf{\kn{ವಿಶ್ವವಿದ್ಯಾಲಯ}}, \emph{viśvavidyālaya}.) A degree? (\textbf{\kn{ಪದವಿ}}, \emph{padavi}.) A lecture? (\textbf{\kn{ಉಪನ್ಯಾಸ}}, \emph{upanyāsa}.) A professor? (\textbf{\kn{ಪ್ರಾಧ್ಯಾಪಕ}}, \emph{prādhyāpaka}.) To instruct? (\textbf{\kn{ಬೋಧಿಸು}}, \emph{bōdhisu}.) Research? (\textbf{\kn{ಸಂಶೋಧನೆ}}, \emph{saṃśōdhane}.) An experiment? (\textbf{\kn{ಪ್ರಯೋಗ}}, \emph{prayōga}.) A theory? (\textbf{\kn{ಸಿದ್ಧಾಂತ}}, \emph{siddhānta}.) To analyse? (\textbf{\kn{ವಿಶ್ಲೇಷಿಸು}}, \emph{viślēṣisu}.) A survey? (\textbf{\kn{ಸಮೀಕ್ಷೆ}}, \emph{samīkṣe}.) To focus? (\textbf{\kn{ಕೇಂದ್ರೀಕರಿಸು}}, \emph{kēndrīkarisu}.) A study? (\textbf{\kn{ಅಧ್ಯಯನ}}, \emph{adhyayana}.) Optional? (\textbf{\kn{ಐಚ್ಛಿಕ}}, \emph{aicchika}.) A syllabus? (\textbf{\kn{ಪಠ್ಯಕ್ರಮ}}, \emph{paṭhyakrama}.) Revision? (\textbf{\kn{ಪುನರಾವರ್ತನೆ}}, \emph{punarāvartane}.) An essay? (\textbf{\kn{ಪ್ರಬಂಧ}}, \emph{prabandha}.) To summarise? (\textbf{\kn{ಸಂಕ್ಷೇಪಿಸು}}, \emph{saṅkṣēpisu}.) A treatise? (\textbf{\kn{ಗ್ರಂಥ}}, \emph{grantha}.) To present? (\textbf{\kn{ಮಂಡಿಸು}}, \emph{maṇḍisu}.) A note? (\textbf{\kn{ಟಿಪ್ಪಣಿ}}, \emph{ṭippaṇi}.) Thus? (\textbf{\kn{ಹೀಗಾಗಿ}}, \emph{hīgāgi}.) However? (\textbf{\kn{ಆದಾಗ್ಯೂ}}, \emph{ādāgyū}.) On account of? (\textbf{\kn{ನಿಮಿತ್ತ}}, \emph{nimitta}.) Above all? (\textbf{\kn{ಮೇಲಾಗಿ}}, \emph{mēlāgi}.) Not only that? (\textbf{\kn{ಮಾತ್ರವಲ್ಲದೆ}}, \emph{mātravallade}.) A condition? (\textbf{\kn{ಷರತ್ತು}}, \emph{ṣarattu}.) An agreement? (\textbf{\kn{ಒಪ್ಪಂದ}}, \emph{oppanda}.) A pact? (\textbf{\kn{ಒಡಂಬಡಿಕೆ}}, \emph{oḍambaḍike}.) To fulfil? (\textbf{\kn{ಈಡೇರಿಸು}}, \emph{īḍērisu}.) An obligation? (\textbf{\kn{ಬಾಧ್ಯತೆ}}, \emph{bādhyate}.)
\end{tcolorbox}
